\subsection{换基域}

我们沿用上一小节的记号。设 \(K'\) 是一个代数闭体，使得元素 \(n \cdot 1\) 在 \(K'\) 中非零。群 \(\mu_{n}(\mathrm{K})\) 与 \(\mu_{n}(\mathrm{K}')\) 都是 n 阶循环群（V, §11, No.~2, 第 78 页，定理 1）。取定一个同构 \(\varphi\)（从 \(\mu_{n}(\mathrm{K})\) 到 \(\mu_{n}(\mathrm{K}')\) 的同构）。设 \(\pi\) 是 G 在有限维 K-向量空间中的一个线性表示，并设 \(\pi'\) 是 G 在有限维 \(K'\)-向量空间中的一个线性表示。我们称 \(\pi\) 与 \(\pi'\)（通过 \(\varphi\)）相关联，如果对每个 \(g \in G\) 与每个 \(\omega \in \mu_{n}(\mathrm{K})\)，\(\omega\) 作为 \(\pi(g)\) 的特征值的重数等于 \(\varphi(\omega)\) 作为 \(\pi'(g)\) 的特征值的重数。这时 \(\pi\) 与 \(\pi'\) 有相同的维数，取 g = 1 即可看出。

设 \(\pi_1\) 与 \(\pi_2\)（相应地，\(\pi_1'\) 与 \(\pi_2'\)）是 \(G\) 在 \(K\)（相应地，\(K'\)）上有限维向量空间中的线性表示。我们有下列性质：

\begin{enumerate}
\def\labelenumi{\alph{enumi})}
\item
  若 \(\pi_1\) 与 \(\pi_1'\) 和 \(\pi_2'\) 都相关联，则 \(\pi_1'\) 与 \(\pi_2'\) 同构。
\item
  若 \(\pi_1\) 与 \(\pi_1'\) 相关联且 \(\pi_2\) 与 \(\pi_2'\) 相关联，则 \(\pi_1 \oplus \pi_2\) 与 \(\pi_1' \oplus \pi_2'\) 相关联，且 \(\pi_1 \otimes \pi_2\) 与 \(\pi_1' \otimes \pi_2'\) 相关联。
\end{enumerate}

断言 a) 由 VIII, 第 402 页的推论 5 得出，而断言 b) 是显然的。

这里，\(\mathcal{S}_{\mathrm{K}}(\mathrm{G})\) 表示单 K{[}G{]}-模的类所成的集合（先前记作 \(\widehat{G}\)）；同样定义 \(\mathcal{S}_{\mathrm{K}^{\prime}}(\mathrm{G})\)。集合 \(\mathcal{S}_{\mathrm{K}}(\mathrm{G})\) 与 \(\mathcal{S}_{\mathrm{K}^{\prime}}(\mathrm{G})\) 都是有限的，其基数为共轭类的个数（VIII, 第 411 页，命题 5）。

命题 10. —— 存在唯一的映射 \(\varphi_{\mathrm{G}}\)，从 \(\mathcal{S}_{\mathrm{K}}(\mathrm{G})\) 到 \(\mathcal{S}_{\mathrm{K}'}(\mathrm{G})\)，使得 \(\lambda\) 与 \(\varphi_{\mathrm{G}}(\lambda)\) 通过 \(\varphi\) 相关联——这里 \(\lambda\) 取遍 \(\mathcal{S}_{\mathrm{K}}(\mathrm{G})\)。此外，\(\varphi_{\mathrm{G}}\) 是双射。

\(\varphi_{G}\) 的唯一性由上面的性质 a) 得出。

\textbf{A) 假设体 K 的特征为 0。}

群 \(\mu_n(\mathrm{K})\) 是循环群（V, §11, No.~2, 第 78 页，定理 1）；取该群的一个生成元 \(\zeta\)。考虑环同态 \(\rho: \mathbf{Z}[\mathrm{X}] \to \mathcal{O}_n\)，它把 X 映为 \(\zeta\)。它是满射。分圆多项式 \(\Phi_n(\mathrm{X})\) 在 \(\mathbf{Q}[\mathrm{X}]\) 中不可约（V, §11, No.~5, 第 84 页，定理 2）；因此它是 \(\zeta\) 在 \(\mathbf{Q}\) 上的最小多项式。多项式 \(\Phi_n\) 是首一的且具有整系数（V, §11, No.~4, 第 81 页）。设 \(\mathrm{P} \in \mathbf{Z}[\mathrm{X}]\) 是满足 \(\mathrm{P}(\zeta) = 0\) 的一个多项式；由多项式的欧几里得除法（IV, §1, No.~6, 第 10 页），存在 \(\mathbf{Z}[\mathrm{X}]\) 中的两个多项式 Q 与 R，使得 \(\mathrm{P} = \mathrm{Q}\Phi_n + \mathrm{R}\) 且 \(\deg(\mathrm{R}) < \deg(\Phi_n)\)。我们有 \(\mathrm{R}(\zeta) = 0\)，从而 \(\mathrm{R} = 0\)，因为 \(\Phi_n\) 是 \(\zeta\) 的最小多项式。于是 \(\rho\) 的核是理想 \(\Phi_n\mathbf{Z}[\mathrm{X}]\)（\(\mathbf{Z}[\mathrm{X}]\) 中的理想），而 \(\rho\) 诱导出从 \(\mathbf{Z}[\mathrm{X}] / \Phi_n\mathbf{Z}[\mathrm{X}]\) 到 \(\mathcal{O}_n\) 的环同构。

令 \(\zeta' = \varphi(\zeta)\)；它是一个 \(n\) 次本原单位根，位于 \(\mathbf{K}'\) 中，因此我们有 \(\Phi_n(\zeta') = 0\)（V, §11, No.~5, 第 83 页，引理 3）。从而存在同态 \(\varphi_0\)（从环 \(\mathcal{O}_n\) 到体 \(\mathbf{K}'\) 的同态），它把 \(\zeta\) 变为 \(\zeta'\)；它是映射 \(\varphi\)（从 \(\mu_n(\mathrm{K})\) 到 \(\mu_n(\mathrm{K}')\) 的映射）的扩张。设 \(\mathcal{O}\) 是 \(\mathrm{K}\) 中由形如 \(\frac{a}{n^r}\)（\(a \in \mathcal{O}_n\)，\(r \in \mathbf{N}\)）的元素组成的子环。由于 \(n \cdot 1\) 在 \(\mathbf{K}'\) 中可逆，同态 \(\varphi_0\) 可扩张为一个同态 \(\varphi_1\)（从 \(\mathcal{O}\) 到 \(\mathbf{K}'\) 的同态）。

我们把 \(\mathcal{O}[\mathrm{G}]\)（群 \(\mathrm{G}\) 在 \(\mathcal{O}\) 上的代数）与代数 \(\mathrm{K}[\mathrm{G}]\) 的一个子环等同起来，并借助下述公式定义环同态 \(\Phi\)（从 \(\mathcal{O}[\mathrm{G}]\) 到 \(\mathrm{K}'[\mathrm{G}]\) 的环同态）：

\[
\Phi \bigg (\sum_ {g \in \mathrm{G}} a _ {g} g \bigg) = \sum_ {g \in \mathrm{G}} \varphi_ {1} (a _ {g}) g.\tag{42}
\]

用 \(\mathcal{C}\) 表示 G 的共轭类所成的集合。对 \(\mathcal{C}\) 中的 C，用 \(u_{\mathrm{C}}\) 表示元素 \(\sum_{g\in \mathrm{C}}g\)（\(\mathcal{O}[\mathrm{G}]\) 的元素）；族 \((u_{\mathrm{C}})_{\mathrm{C}\in \mathcal{C}}\) 构成 \(\mathcal{O}\) 上向量空间 \(\mathrm{Z}(\mathcal{O}[\mathrm{G}])\)（代数 \(\mathcal{O}[\mathrm{G}]\) 的中心）的一组基（VIII, 第 398 页）。对任意元素 \(\lambda\)（\(\mathcal{S}_{\mathrm{K}}(\mathrm{G})\) 的元素），用 \(\chi_{\lambda}\) 表示它的特征标，用 \(d_{\lambda}\) 表示它的维数，并用 \(e_{\lambda}\) 表示

由下式定义的 K{[}G{]} 中的元素

\[
e _ {\lambda} = | \mathrm{G} | ^ {- 1} d _ {\lambda} \sum_ {g \in \mathrm{G}} \chi_ {\lambda} (g ^ {- 1}) g.\tag{43}
\]

族 \((e_{\lambda})\) 是 \(Z(K[G])\)（环 K{[}G{]} 的中心）在 K 上的一组基（VIII, 第 408 页）。我们有

\[
e _ {\lambda} = \sum_ {\mathrm{C} \in \mathcal {C}} \alpha_ {\lambda , \mathrm{C}} u _ {\mathrm{C}}\tag{44}
\]

其中

\[
\alpha_ {\lambda , \mathrm{C}} = | \mathrm{G} | ^ {- 1} d _ {\lambda} \chi_ {\lambda} (\mathrm{C} ^ {- 1}).\tag{45}
\]

对 \(C \in \mathcal{C}\) 与 \(\lambda \in \mathcal{S}_{\mathrm{K}}(\mathrm{G})\)，令 \(\beta_{\mathrm{C},\lambda} = |\mathrm{G}| d_{\lambda}^{-1} d(\mathrm{C})^{-1} \chi_{\lambda}(\mathrm{C})\)。矩阵 \((\alpha_{\lambda,\mathrm{C}})\) 的元素都属于 \(\mathcal{O}\)，并且由 VIII, 第 411 页的公式 (31)，它的逆矩阵是矩阵 \((\beta_{\mathrm{C},\lambda})\)，而后者的元素由 VIII, 第 415 页的命题 9 也都属于 \(\mathcal{O}\)。于是族 \((e_{\lambda})\) 是 \(\mathcal{O}\)-模 \(\mathrm{Z}(\mathcal{O}[\mathrm{G}])\) 的一组基。

诸元素 \(\Phi(u_{\mathrm{C}}) = \sum_{g\in \mathrm{C}}g\)（\(K^{\prime}[G]\) 中的元素）构成 \(K^{\prime}\) 上向量空间 \(Z(K^{\prime}[G])\)（环 \(K^{\prime}[G]\) 的中心）的一组基。我们有

\[
\Phi (e _ {\lambda}) = \sum_ {\mathrm{C} \in \mathcal {C}} \varphi_ {1} (\alpha_ {\lambda , \mathrm{C}}) \Phi (u _ {\mathrm{C}}),\tag{46}
\]

并且以 \(\varphi_{1}(\alpha_{\lambda,\mathrm{C}})\) 为元素的矩阵可逆。因此诸 \(\Phi(e_{\lambda})\) 所成的族是 \(\mathrm{Z}(\mathrm{K}^{\prime}[\mathrm{G}])\) 的一组基。族 \((e_{\lambda})\) 是 \(\mathrm{Z}(\mathrm{K}[\mathrm{G}])\) 中幂等元 1 的一个分拆（VIII, 第 146 页与第 408 页）；换言之，我们有

\[
\sum_ {\lambda} e _ {\lambda} = 1, \quad e _ {\lambda} ^ {2} = e _ {\lambda}, \quad e _ {\lambda} e _ {\mu} = 0 \quad \mathrm{if} \quad \lambda \neq \mu .
\]

由此推出，诸 \(\Phi(e_{\lambda})\) 所成的族是 \(Z(K^{\prime}[G])\) 中幂等元 1 的一个分拆；由于该族是 \(Z(K^{\prime}[G])\) 在 \(K^{\prime}\) 上的一组基，其元素是 \(Z(K^{\prime}[G])\) 中的不可分解幂等元（VIII, 第 148 页，注 4）。

对 \(\lambda'\) 属于 \(\mathcal{S}_{\mathrm{K}'}(\mathrm{G})\) 的情形，如上定义 \(\chi_{\lambda'}, d_{\lambda'}\) 与 \(e_{\lambda'}\)。由 VIII, 第 408 页，诸元素 \(e_{\lambda'}\) 是 \(\mathrm{Z}(\mathrm{K}'[\mathrm{G}])\) 中的不可分解幂等元。因此存在双射 \(\varphi_{\mathrm{G}}\)，从 \(\mathcal{S}_{\mathrm{K}}(\mathrm{G})\) 到 \(\mathcal{S}_{\mathrm{K}'}(\mathrm{G})\)，使得 \(\Phi(e_{\lambda}) = e_{\varphi_{\mathrm{G}}(\lambda)}\)，其中 \(\lambda\) 取遍 \(\mathcal{S}_{\mathrm{K}}(\mathrm{G})\)。

设 \(\lambda\) 是 \(\mathcal{S}_{\mathrm{K}}(\mathrm{G})\) 中的元素；令 \(\lambda' = \varphi_{\mathrm{G}}(\lambda)\)。设 \((\mathrm{V}_{\lambda}, \pi_{\lambda})\)（相应地，\((\mathrm{V}_{\lambda'}, \pi_{\lambda'})\)）是 G 的一个线性表示，其相伴的 K{[}G{]}-模（相应地，K'{[}G{]}-模）的类为 \(\lambda\)（相应地，\(\lambda'\)）。我们来证明 \(\lambda\) 与 \(\lambda'\) 相关联。设 \(g\) 是 G 的一个元素。设 \(\delta(\mathrm{T})\) 是自同态 \(1 + \mathrm{T}\pi_{\lambda}(g)\)

（K{[}T{]}-模 K{[}T{]} \(\otimes_{\mathrm{K}}\mathrm{V}_{\lambda}\) 的自同态）的行列式。用 \(\omega_{1},\ldots ,\omega_{d_{\lambda}}\) 表示 \(\pi_{\lambda}(g)\) 的特征值；我们有

\[
\delta (\mathrm{T}) = (1 + \mathrm{T} \omega_ {1}) \dots (1 + \mathrm{T} \omega_ {d _ {\lambda}}).
\]

我们同样定义 \(\delta'(T)\)，并把 \(\pi_{\lambda'}(g)\) 的特征值记作 \(\omega_1', \ldots, \omega_{d_{\lambda'}}'\)。\(\mathcal{O}\)-模 \(\mathcal{O}[G]\) 是自由的，以 \(G\) 为基，而 \(K\)-向量空间 \(K[G]\) 以 \(G\) 为基。用 \(\Delta(T)\) 表示乘以 \(1 + e_{\lambda}gT\)（在 \(\mathcal{O}[T]\)-模 \(\mathcal{O}[T] \otimes_{\mathcal{O}} \mathcal{O}[G]\) 中）的行列式。它也是乘以 \(1 + e_{\lambda}gT\)（在 \(K[T]\)-模 \(K[T] \otimes_K K[G]\) 中）的行列式。设 \(\overline{\varphi}_1\) 是从 \(\mathcal{O}[T]\) 到 \(K'[T]\) 的同态，它扩张 \(\varphi_1\) 且把 \(T\) 映为 \(T\)。由于 \(G\) 是 \(K'\)-向量空间 \(K'[G]\) 的一组基，多项式 \(\overline{\varphi}_1(\Delta(T))\) 等于行列式 \(\Delta'(T)\)，即乘以 \(1 + e_{\lambda'}gT\)（在 \(K'\)-向量空间 \(K'[G]\) 中）的行列式。

代数 K{[}G{]} 是它的诸单分量 \(e_{\mu}K[G]\)（\(\mu\) 取遍 \(\mathcal{S}_{\mathrm{K}}(\mathrm{G})\)）的直和。当 \(\mu\) 异于 \(\lambda\) 时，元素 \(e_{\lambda}g\) 零化 \(e_{\mu}K[G]\)。此外，乘以 \(e_{\lambda}g\) 的乘法与乘以 g 的乘法在 \(e_{\lambda}K[G]\) 中一致。注意到 VIII, 第 409 页以及 VIII, 第 400 页的例 6，G 在 \(e_{\lambda}K[G]\) 中的表示是 \(d_{\lambda}\) 个类为 \(\lambda\) 的表示的直和。因此我们有 \(\Delta(\mathrm{T}) = \delta(\mathrm{T})^{d_{\lambda}}\)。

同理，我们有 \(\Delta^{\prime}(\mathrm{T}) = \delta^{\prime}(\mathrm{T})^{d_{\lambda^{\prime}}}\)。

由关系式 \(\Delta'(\mathrm{T}) = \overline{\varphi}_1(\Delta(\mathrm{T}))\)，我们首先推出 \(d_{\lambda}^2 = d_{\lambda'}^2\)，从而 \(d_{\lambda} = d_{\lambda'}\)；然后推出序列 \(\varphi(\omega_1), \ldots, \varphi(\omega_{d_\lambda})\) 可由序列 \((\omega_1', \ldots, \omega_{d_{\lambda'}'}')\) 经过指标集的一个置换得到。

由于这一点对每个元素 \(g\)（\(G\) 的元素）都成立，表示 \(\lambda\) 与 \(\lambda'\) 相关联。这样，当体 \(K\) 的特征为 0 时，我们就证明了命题 10。

\textbf{B) 一般情形。}

设 L 是一个特征为 0 的代数闭体（例如 Q 的一个代数闭包）。用 \(\mathcal{S}_{\mathrm{L}}(\mathrm{G})\) 表示单 L{[}G{]}-模的类所成的集合。取定一个同构 \(\eta\)（从群 \(\mu_{n}(\mathrm{L})\) 到群 \(\mu_{n}(\mathrm{K})\) 的同构），并令 \(\eta' = \varphi \circ \eta\)。由证明的部分 A)，存在双射

\[
\eta_ {\mathrm{G}}: \mathcal {S} _ {\mathrm{L}} (\mathrm{G}) \to \mathcal {S} _ {\mathrm{K}} (\mathrm{G}), \quad \eta_ {\mathrm{G}} ^ {\prime}: \mathcal {S} _ {\mathrm{L}} (\mathrm{G}) \to \mathcal {S} _ {\mathrm {K^ {\prime}}} (\mathrm{G})
\]

它们具有下述性质：对每个 \(\lambda\)（\(\mathcal{S}_{\mathrm{L}}(\mathrm{G})\) 中的元素），表示 \(\lambda\) 与 \(\eta_{\mathrm{G}}(\lambda)\) 通过 \(\eta\) 相关联，且表示 \(\lambda\) 与 \(\eta_{\mathrm{G}}^{\prime}(\lambda)\) 通过 \(\eta^{\prime}\) 相关联。双射 \(\varphi_{G} = \eta_{G}^{\prime} \circ \eta_{G}^{-1}\) 具有所要的性质。

双射 \(\varphi_{G}\)（从 \(\mathcal{S}_{K}(G)\) 到 \(\mathcal{S}_{K'}(G)\) 的双射）可扩张为一个同构（仍记作 \(\varphi_{G}\)），即从格罗滕迪克群 \(R_{K}(G)\) 到群 \(R_{K'}(G)\) 的同构。

注 1. —— 假设 \(K'\) 是 K 的一个扩张，且同构 \(\varphi\) 是映射 \(\xi \mapsto \xi \cdot 1\)；那么映射 \(\varphi_{G}\) 由从 K 到 \(K'\) 的标量扩张给出。

推论 1. —— 映射 \(\varphi_{\mathrm{G}}\) 是从 \(\mathrm{R}_{\mathrm{K}}(\mathrm{G})\) 到 \(\mathrm{R}_{\mathrm{K}^{\prime}}(\mathrm{G})\) 的环同构。对每个有限维表示 \(\pi\)（\(\mathrm{G}\) 在 K-向量空间中的表示），我们有 \(\varphi_{\mathrm{G}}([\pi]) = [\pi^{\prime}]\)，其中 \(\pi^{\prime}\) 是与 \(\pi\) 通过 \(\varphi\) 相关联的一个表示。

这由 G 的表示的半单性以及 VIII, 第 416 页的性质 b) 得出。

推论 2. —— G 的每个单表示的维数整除 G 的阶。

这由命题 10 以及 VIII, 第 415 页的命题 9 得出。

注. —— 2) 假设群 G 是交换群。我们在 VIII, 第 414 页的注中已看到，\(\mathcal{S}_{\mathrm{K}}(\mathrm{G})\) 可与集合 \(\operatorname{Hom}(\mathrm{G}, \mu_n(\mathrm{K}))\) 等同。同样，\(\mathcal{S}_{\mathrm{K}'}(\mathrm{G})\) 可与 \(\operatorname{Hom}(\mathrm{G}, \mu_n(\mathrm{K}'))\) 等同。在这些等同之下，双射 \(\varphi_{\mathrm{G}}\) 就是映射 \(\chi \mapsto \varphi \circ \chi\)。

\begin{enumerate}
\def\labelenumi{\arabic{enumi})}
\setcounter{enumi}{2}
\tightlist
\item
  设 \(\pi_1\) 与 \(\pi_2\) 是 G 在 K 上有限维向量空间中的线性表示。对 \(i = 1,2\)，设 \(\pi_i'\) 是与 \(\pi_i\) 通过 \(\varphi\) 相关联的一个表示。我们有
\end{enumerate}

\[
\dim_ {K} \operatorname{Hom} _ {K} (\pi_ {1}, \pi_ {2}) = \dim_ {K ^ {\prime}} \operatorname{Hom} _ {K ^ {\prime}} (\pi_ {1} ^ {\prime}, \pi_ {2} ^ {\prime}).\tag{47}
\]

证明沿用 VIII, 第 410 页的推论的证明，归结为诸 \(\pi_{i}\)（从而诸 \(\pi_{i}^{\prime}\)）都是单表示的情形。

\begin{enumerate}
\def\labelenumi{\arabic{enumi})}
\setcounter{enumi}{3}
\tightlist
\item
  设 H 是 G 的一个基数为 m 的子群。同构 \(\varphi\) 限制为从 \(\mu_{m}(\mathrm{K})\) 到 \(\mu_{m}(\mathrm{K}^{\prime})\) 的同构，从而给出一个环同构 \(\varphi_{H}\)，从 \(\mathrm{R}_{\mathrm{K}}(\mathrm{H})\) 到 \(\mathrm{R}_{\mathrm{K}^{\prime}}(\mathrm{H})\)。下面的图交换：
\end{enumerate}

\[
\begin{array}{c c} \mathrm {R_ {K} (G)} \xrightarrow {\mathrm {Res_ {H} ^ {G}}} \mathrm {R_ {K} (H)} & \qquad \qquad \mathrm {R_ {K} (H)} \xrightarrow {\mathrm {Ind_ {H} ^ {G}}} \mathrm {R_ {K} (G)} \\ \Biggl \downarrow \varphi_ {\mathrm{G}} & \Biggl \downarrow \varphi_ {\mathrm{H}} \\ \mathrm {R_ {K^ {\prime}} (G)} \xrightarrow {\mathrm {Res_ {H} ^ {G}}} \mathrm {R_ {K^ {\prime}} (H)}, & \qquad \qquad \mathrm {R_ {K^ {\prime}} (H)} \xrightarrow {\mathrm {Ind_ {H} ^ {G}}} \mathrm {R_ {K^ {\prime}} (G)}. \end{array}
\]

第一个图的交换性是显然的，而第二个图的交换性由它借助弗罗贝尼乌斯互反律与公式 (47) 得出。

\begin{enumerate}
\def\labelenumi{\arabic{enumi})}
\setcounter{enumi}{4}
\tightlist
\item
  假设 \(G\) 是两个有限群的乘积 \(G' \times G''\)。我们按上例的方式定义同构 \(\varphi_{G'}\) 与 \(\varphi_{G''}\)。
\end{enumerate}

我们有一个交换图

\[
\begin{array}{c} \mathrm{R} _ {\mathrm{K}} (\mathrm{G} ^ {\prime}) \otimes_ {\mathbf {Z}} \mathrm{R} _ {\mathrm{K}} (\mathrm{G} ^ {\prime \prime}) \xrightarrow {\kappa} \mathrm{R} _ {\mathrm{K}} (\mathrm{G}) \\ \Biggl \downarrow \varphi_ {\mathrm{G} ^ {\prime}} \otimes \varphi_ {\mathrm{G} ^ {\prime \prime}} \qquad \qquad \qquad \Biggl \downarrow \varphi_ {\mathrm{G}} \\ \mathrm{R} _ {\mathrm{K} ^ {\prime}} (\mathrm{G} ^ {\prime}) \otimes_ {\mathbf {Z}} \mathrm{R} _ {\mathrm{K} ^ {\prime}} (\mathrm{G} ^ {\prime \prime}) \xrightarrow {\kappa^ {\prime}} \mathrm{R} _ {\mathrm{K} ^ {\prime}} (\mathrm{G})  , \end{array}
\]

其中同构 \(\kappa\) 与 \(\kappa'\) 是 VIII, 第 406 页中定义的那两个同构。

